\section{Robustness}
\label{sec:robustness}

Three assumptions carry more weight than the rest, and each is worth a paragraph
that a reader can disagree with directly.

\paragraph{Depreciation.} Setting $\delta = 0.06$ is conventional rather than
estimated. Re-running the perpetual inventory recursion~\eqref{eq:pim} at
$\delta = 0.04$ and $\delta = 0.08$ shifts the level of the capital series
noticeably but leaves $\Delta \log K$ nearly untouched, so $\hat{\alpha}$ moves
by less than a point. The capital-output ratio in Table~\ref{tab:summary} is not
so lucky: it ranges from $2.4$ to $3.4$ across the three choices, and since that
moment is what identifies $\alpha$ in the calibrated version of the model, the
two routes to $\alpha$ disagree by more than either one's standard error admits.

\paragraph{Constant returns.} Imposing $\alpha + (1 - \alpha) = 1$
in~\eqref{eq:production} is what allows the factor shares in
Table~\ref{tab:estimates} to be read off a single coefficient. Relaxing it and
estimating the two elasticities freely gives $0.36$ and $0.61$, whose sum is
statistically indistinguishable from one and economically indistinguishable from
anything else in the neighbourhood. This is reassuring without being evidence.

\subsection{Measurement of the residual}

Because $\Delta \log A_{t}$ is defined as what~\eqref{eq:growth} does not
explain, every measurement error in output, capital or employment lands in it
with the wrong sign attached. The persistence of the estimated residual is
therefore an upper bound on the persistence of true productivity, and the
variance decomposition in Section~\ref{sec:data} should be read with that in
mind.

None of this changes the qualitative picture, which is the sentence every
robustness section ends on and the one readers most enjoy clicking to complain
about. What it does change is how much confidence the twenty per cent band
around $K^{\star}$ deserves, and the honest answer is: less than the tables
imply, more than a sceptic would grant.
